% Corrected French for printed p.202.
% Thirteen print-level repairs are recorded append-only.
\clearpage
\SourcePage{202}

\LayerTarget{tohoku:proof:proposition-5.2.2:p1}
\noindent foncteur de ${\mathbf C}^{X(G)}$ dans ${\mathbf C}^Y$ ; mais en vertu du calcul de ces
foncteurs dérivés par le lemme 3.7.2, et se rappelant qu'un objet injectif de
${\mathbf C}^{X(G)}$ est un faisceau flasque, on voit qu'on trouve le même résultat que
si on considère $f_*$ comme un foncteur de ${\mathbf C}^X$ dans ${\mathbf C}^Y$, i.e.
\textit{$R^qf_*(A)$ est le faisceau sur $Y$ associé au préfaisceau formé des
$H^q(f^{-1}(U),A)$.}

\LayerTarget{tohoku:remark:cohomological-hypotheses-5.2}
L'hypothèse de la proposition 5.2.2 doit être regardée comme un équivalent
cohomologique de l'hypothèse que $G$ est un ``groupe discontinu d'opérateurs''
sur $X$. Un autre cas important est celui où $\mathbf H^q(G,A)=0$ pour $q>0$,
hypothèse qui doit être regardée comme un équivalent cohomologique de la
condition usuelle que ``$G$ opère sans points fixes'' :

\LayerTarget{tohoku:proposition:5.2.3}
\textsc{Proposition 5.2.3.} \textit{Supposons que
$\mathbf H^q(G,A)=0$ pour $q>0$. Alors $H^n(X;G,A)$ est isomorphe à
$H^n(Y,A^G)$, donc $H^*(Y,A^G)$ est l'aboutissement d'une suite spectrale
cohomologique dont le terme initial est
$II_2^{p,q}(A)=H^p(G,H^q(X,A))$.}

\LayerTarget{tohoku:proof:proposition-5.2.3}
C'est une conséquence immédiate de la première suite spectrale du théorème
5.2.1.
\LayerTarget{tohoku:application:classical-spaces-with-operators}
On obtient ainsi la forme classique de la théorie des espaces à opérateurs
[4]. Nous donnerons au numéro suivant des conditions de validité de cette
proposition, $G$ étant un ``groupe discontinu sans points fixes d'opérateurs
dans $X$''. D'autres conditions de validité sont liées à la caractéristique
d'un corps de base $k$ pour le faisceau $A$ et les ordres des stabilisateurs
$G_y$ ; un cas particulièrement simple est le suivant (qui se démontre
d'ailleurs directement très facilement, et est sans doute bien connu) :

\LayerTarget{tohoku:corollary:5.2.3}
\textsc{Corollaire.} \textit{Supposons outre $\mathbf H^q(G,A)=0$
pour $q>0$ que $G$ soit un groupe fini d'ordre $m$, que l'espace $X$ soit
séparé, et que la multiplication par $m$ dans $A$ soit un automorphisme de
$A$ (par exemple $A$ est un faisceau d'espaces vectoriels sur un corps de
caractéristique ne divisant pas $m$). Alors on a
$H^n(Y,A^G)=H^n(X,A)^G$, les deux membres étant isomorphes à
$H^n(X;G,A)$.}

\LayerTarget{tohoku:transition:final-particular-case-5.2}
Un dernier cas particulier intéressant est le suivant :

\LayerTarget{tohoku:proposition:5.2.4}
\textsc{Proposition 5.2.4.} \textit{Supposons $H^q(X;A)=0$ pour
$q>0$, alors $H^n(X;G,A)=H^n(G,\Gamma(A))$, et par suite
$H^*(G,\Gamma(A))$ est l'aboutissement d'une suite spectrale dont le terme
initial est $I_2^{p,q}(A)=H^p(Y,\mathbf H^q(G,A))$.}

\LayerTarget{tohoku:application:modular-group-cohomology}
Cette proposition peut être utilisée pour le calcul de la cohomologie de
certains groupes, par exemple la cohomologie à coéfficients constants de
divers groupes modulaires à une variable [10].

\LayerTarget{tohoku:remark:5.2.4:1}
\textsc{Remarques. 1.} Supposons qu'on se donne un $G$-faisceau
${\mathbf O}$, et qu'on se borne à considérer des $G$-${\mathbf O}$-Modules, $\Gamma_X^G$ étant
donc considéré comme un foncteur de ${\mathbf C}^{\mathrm O(G)}$ dans la catégorie des modules
sur $\Gamma({\mathbf O})^G$. Alors ses foncteurs dérivés coïncident avec les foncteurs
$H^p(X;G,A)$ précédents, comme il résulte aussitôt du fait que tout objet
injectif de ${\mathbf C}^{\mathrm O(G)}$ est annulé par les $H^n(X;G,A)$ pour $n>0$. Ce dernier
fait sera prouvé dans le lemme 5.6.1, plus bas (où on remplace $B$ par $A$, $A$
par ${\mathbf O}$).

\LayerTarget{tohoku:remark:5.2.4:2}
2. On peut, des deux suites exactes (5.2.8) et (5.2.9), éliminer
$H^*(X;G,A)$ de façon à obtenir des relations entre la cohomologie de $Y$,
celle de $X$

\TohokuContinuation{203}
